\documentclass{jsarticle}
\usepackage[dvipdfmx, final]{graphicx}
\usepackage{amsmath}
\usepackage{amssymb}
\usepackage{chemfig}
\begin{document}
\title{Euler product estrade from \\
Heisenberg Non-commutative with \\
deprivate equation}
\author{Masaaki Yamaguchi}
\date{}
\maketitle

   \newcommand{\variint}{%
	\begin{picture}(15,30)
		\put(0,0){
			$ \iint $}
		\put(0,5){\line(15,0){15}}
	\end{picture} %
}
   \newcommand{\vare}{%
	\begin{picture}(5,5)
		\put(4,4){
			 \line(1,1){5} }
		\put(0,0){$\vee$}
	\end{picture} %
}
   \newcommand{\varve}{%
	\begin{picture}(5,5)
		\put(0,0){$\vee$}
		\put(0,5){\line(5,0){5}}
	\end{picture} %
}
\newcommand{\variiint}{%
	\begin{picture}(15,15)
		\put(0,0){
			$ \iiint $}
		\put(0,5){\line(15,0){20}}
	\end{picture} %
}



 \newcommand{\alearletter}{%
	 \begin{picture}(10,10)
		 \put(0,0){
			 $ \square $}
		 \put(0,0){\line(1,1){10}}
	 \end{picture} %
	 }

 \newcommand{\secproduct}{%
	 \begin{picture}(10,20)
		 \put(0,0){
			 $ \nabla $}
		 \put(5,0){+}
	 \end{picture} %
	 }
 \newcommand{\talot}{%
	 \begin{picture}(15,20)
		 \put(0,0){
			 $ \text{\scalebox{2.0}[2]{$\nabla$}}$}
		 \put(5,3){
			 $ \Delta $}
	 \end{picture} %
	 }
 \newcommand{\hazerdnabla}{%
	 \begin{picture}(15,20)
		 \put(0,0){
			 $ \text{\scalebox{2.0}[2]{$\nabla$}}$}
		 \put(5,3){
			 $\text{Y}  $}
	 \end{picture} %
	 }

 \newcommand{\securebox}{%
	 \begin{picture}(15,20)
		 \put(0,0){
			 $ \text{\scalebox{2.0}[2]{$\square$}}$}
		 \put(5,3){
			 $\text{Y}  $}
	 \end{picture} %
	 }


\newcommand{\timebow}{%
	\begin{picture}(5,5)
		\put(0,0){$\text{\rotatebox{90}{$\bowtie$}}$}
	\end{picture} %
	}


 \newcommand{\flowerletter}{%
	 \begin{picture}(10,20)
		 \put(0,0){$\text{\scalebox{1.0}[1]{$\square$}}$}
		 \put(0,0){$\text{\scalebox{1.0}[2]{$\int$}}$}
	 \end{picture} %
	 }
\newcommand{\squareiint}{%
	\begin{picture}(15,30)
		\put(0,0){
			$ \iint $}
		\put(0,5){$\text{\scalebox{2.0}[1]{$\square$}}$}
	\end{picture} %
}

\newcommand{\squareint}{%
	\begin{picture}(15,30)
		\put(0,0){
			$ \int $}
		\put(0,5){$\text{\scalebox{2.0}[1]{$\square$}}$}
	\end{picture} %
}


   \newcommand{\dazanier}{%
	\begin{picture}(12,12)
		\put(0,0){
			$ \bigsqcup $}
		\put(4,3){$\text{\scalebox{1.0}[1]{$\top$}}$}
		\put(4,3){$\text{\scalebox{1.0}[1]{$\bot$}}$}
	\end{picture} %
}
\newcommand{\innabla}{%
	\begin{picture}(13,13)
		\put(0,0){
			$ \nabla $}
		\put(0,0){$\text{\scalebox{2.0}[2]{$\nabla$}}$}
	\end{picture} %
}

\newcommand{\inndelta}{%
	\begin{picture}(20,20)
		\put(0,0){
			$ \Delta $}
		\put(0,0){$\text{\scalebox{2.0}[2]{$\Delta$}}$}
	\end{picture} %
}

\newcommand{\starnabla}{%
	\begin{picture}(20,20)
		\put(3,5){
			$ \Delta $}
		\put(0,0){$\text{\scalebox{2.0}[2]{$\nabla$}}$}
	\end{picture} %
}


  \newcommand{\whitewithbodercircle}
  {\text{\textcircled{\phantom{\textbullet}}}
  \kern-0.9em\text{\textcircled{\phantom{}}}}
     

 
  \newcommand{\whitewithblack}{\text{\textcircled{\textbullet}}}
  \newcommand{\whitewithblacksquare}{\text{\fbox{\text{\rule{0.5em}
  {0.5em}}}}}
  \newcommand{\whitewithwhitesquare}{\text{\fbox{\text{\phantom
  {\rule{0.5em}{0.5em}}}}}}
 \newcommand{\whitewithbodercircleiiint}{%
	 \begin{picture}(20,20)
		 \put(0,0){\scalebox{2.0}[2]{$\whitewithbodercircle$}}
		 \put(0,0){\scalebox{2.0}[2]{$\iiint$}}
	 \end{picture}%
	 }
\newcommand{\whitewithbodercircleint}{%
	 \begin{picture}(20,20)
		 \put(2,2){\scalebox{1.0}[1]{$\whitewithbodercircle$}}
		 \put(2,2){\scalebox{1.0}[2]{$\int$}}
	 \end{picture}%
	 }
\newcommand{\whitewithbodercircleiint}{%
	 \begin{picture}(20,20)
		 \put(0,0){\scalebox{1.0}[1]{$\whitewithbodercircle$}}
		 \put(2,2){\scalebox{1.0}[2]{$\iint$}}
	 \end{picture}%
	 }
 
 \newcommand{\whitewithbodercirclevarint}{%
	 \begin{picture}(13,13)
		 \put(0,0){\scalebox{2.0}[2]{$\whitewithbodercircle$}}
		 \put(0,0){\scalebox{2.0}[2]{$\varint$}}
	 \end{picture}%
	 }
 
 \newcommand{\whitewithbodercirclevariiint}{%
	 \begin{picture}(13,13)
		 \put(0,0){\scalebox{2.0}[2]{$\whitewithbodercircle$}}
		 \put(0,0){\scalebox{2.0}[2]{$\variiint$}}
	 \end{picture}%
	 }
 
 \newcommand{\average}{%
	  \begin{picture}(30,45)
		  \put(0,0){\scalebox{2.0}[2]{$\square$}}
		  \put(4,4){$\setminus$}
	  \end{picture}%
  } 
  \newcommand{\smallaverage}{%
	  \begin{picture}(30,45)
		  \put(0,0){\scalebox{1.0}[1]{$\square$}}
		  \put(0,0){$\setminus$}
	  \end{picture}%
  } 
 
  \newcommand{\whitewithboder}{%
	  \begin{picture}(15,30)
		  \put(0,0){\scalebox{2.0}[2]{$\square$}}
		  \put(4,4){$\square$}
	  \end{picture}%
  } 
  \newcommand{\whitewithnabla}{%
	  \begin{picture}(15,30)
		  \put(0,0){\scalebox{2.0}[2]{$\nabla$}}
		  \put(3,3){$\nabla$}
	  \end{picture}%
  } 


 Under equations are circle element of step with neipa number represent
 from antigravity equation of global integral manifold quate with
 dalanversian equation, this system neipia number of step with circle
 element also represent with zeta function and squart of g function.
   $$ \int {e^{-x^2 -y^2}}dxdy = \pi $$
   $$ \pi^{e} = (\int e^{\cos \theta + i\sin \theta}d\theta)^{e} $$
   $$ = (\int e^{i\theta}d\theta)^{e} $$
 This system use with Jones manifold and shanon entropy equatioon.
   $$ e^{f} \to f=1, x\log x=1, \int e^{-i\theta}d\theta = \pi^{e} $$
 Antigravity equation also represent circle function.
    $$\alearletter = 2(\sin(ix\log x) + \cos(ix\log x))$$
    Dalanversian equation also represent circle function.
    $$ \square = \cos(ix \log x) - i\sin(ix\log x)$$
    These function recicle with circle of neipa number step function.
    $$ \int e^{-\square}d\square = \pi^{e} $$
    $$ x = {C \over {\log x}}, C = \int{1 \over {x^s}}dx - \log x $$
   Euler product represent with reverse of zeta function. 
   $$ x^y = {1 \over {y^x}}, \pi^{e} = \int e^{-\square}d\square
   = e^{\pi}\int e^{\alearletter}d\alearletter $$
   These equations quaote with being represented with being emerged
   from beta function.
   $$ \square = \alearletter \boxtimes \Psi \to \square = \Psi 
   \boxtimes \alearletter $$
   $$ {d \over {d l}}\square(H\Psi)^{\nabla}, \square{d \over
   {d l}}(H\Psi)^{\nabla} $$
   These function also emerge from global differential equation.
   $$ \beta^{\square -{{\square \over {\log x}}}} =
   \beta^{\square - \alearletter} $$
    $${}^{t}\scalebox{1}[2]{\variiint}
    {\nabla \over {\nabla l}}\square(H\Psi)^{\nabla}
    d\nabla $$
    Varint equation estimate with fundament function. 
    $$ = \int \pi(\chi,\square)^{\nabla}d\nabla $$
    $$ = {}^{\vee}\scalebox{1}[2]{\variint} \pi(\square)d\nabla_m $$
    These system recicle with under environment of equation.
    $$ e^{\pi} = e^{\int e^{-x^2 -y^2}dxdy} = x^y = {1 \over {y^x}}
    = \pi^{e} ~ {1 \over {e^{\pi}}} = {1 \over {\int e^{-\square}d
    \square}} $$
    $$ e^{\pi} = {{\int e^{-\square}d\square} \over {{\int
    e^{\alearletter}d\alearletter}}} $$
    $$ \pi^{e} = e^{\pi} \int e^{\alearletter}d\alearletter $$
    $$ e^{\pi} = {\pi^{e} \over {{\int e^{\alearletter}d\alearletter}}} $$
    $$ \pi^{e} = (\int e^{-(\cos \theta + i\sin \theta)}d\theta)^{e} $$
    $$ \pi^{e} = {1 \over {e^{\pi}}} = {1 \over {{\int e^{-\square}
    d\square}}} $$
    $$ \alearletter = 2(\sin(ix \log x) + \cos(ix \log x)) $$ 
    $$ \square = \cos(ix \log x) - i\sin(ix\log x) $$
    $$ \int e^{-\square}d\square = \pi^{e} $$
    This result with computation escourt with neipa number of step functioon,
    and circle element of pai also zeta function from logment system.
    $$ \log e^{\pi} = \log \left({\pi^{e} \over {{\int e^{\alearletter}
    d\alearletter}}}\right) $$
    This section of equation also represent with pai number comment
    with dalanversian quato with anti-gravity equation.
    $$ \pi = {\square \over {\alearletter}} $$
    And this result also represent with pai number represent with
    beta function of global integral manifold.
    $$ \pi = e^{\beta}\int \beta dx_m $$
    And, spectrum focus of pai number represent with logment function
    into being emerged from g of squart function, and
    this express with function escourt into rout of g conclude with
    zeta function, and universe of formula is pai of circle of reverse
    of formation.
    $$ \pi = {1 \over {\log x}} = \sqrt{g} $$
    $$ = 1 = \pi r^2, r = {1 \over {\sqrt{\pi}}} $$
    After all, this system of neipa and pai number of step function
    represent with colmogoloph function.
    $$ {1 \over {\sqrt{\pi}}} = \int {1 \over {\log x}}dx, 
    \pi = \int \log x dx_m $$
    This equation of facility estimate with fundamantal group of zeta
    function of being reverse of function, this function equal with 
    antigravity and gravity equation from relate of global differential
    manifold escourt into non-communication of equation in pandle of pond
    system.
    $$ \pi(\chi,x) = [i\pi(\chi,x),f(x)] $$
    $$i \pi(\chi,x)\circ f(x) = i\int x\log xdx, f(x)\circ
    \pi(\chi,x) = \int {1 \over {(x\log x)}}dx $$
    This enter facility equal with non-commutative equation.
    $$ \pi(\chi,x) = [i\pi(\chi,x),f(x)] $$
    $$ \pi(\chi,x) = \int x \log xdx $$
    $$ \pi(\chi,x) = \int {1 \over {x \log x}}dx $$
    This of being up equation equal with same energy of entropy in 
    imaginary and reality formula. Anf this system of feed enter with
    shanon entropy in pond of function. Moreover, this energy use
    imaginary pole to equal with integral of shanon entropy and integral
    of colmogloph entropy.
    $$ \pi(\chi,x) = i\pi(chi,x), i\int xlog xdx \cong \int
    {1 \over {(x\log x)}}dx $$
    $$ \int {1 \over {(x\log x)}}dx=i \int x \log xdx - \int{1 \over {
	    (x \log x)}}dx $$
    This equation computative enter into being antigravity equation.
    $$ \int \int {1 \over {(x\log x)^{2}}}dx_m = {1 \over 2}i $$
    This system encall with global integral and differential manifold.
    $$ {d \over {d f}}F(x,y) = {d \over {d f}}\int\int{
	    1 \over {(x\log x)^{2}}}dx_m
    + {d \over {d f}}\int\int{1 \over {(y\log y)^{1 \over 2}}}dy_m $$ 
    This result equation express imaginary number with degree
    of light element of formula.
    $$ i = x^{90^{\circ}}, x\sin 90^{\circ} = i$$
    $$ i = x^{1 \over 2},x= -1 $$
    This imaginary number of reverse is reverse of imaginary result with
    $$ \pi(\chi,x)^{f(x)} = i \int{1 \over {(x\log x)}}\circ f(x)dx $$
    $$ = i \int x \log x dx $$
    $$ {}^{f(x)}\pi(\chi,x) = {}^{f(x)}\int{1 \over {(x\log x)}}dx $$
    $$ = \int {1 \over {(x\log x)}}dx $$
    $$ iy = x\sin 90^{\circ} $$
    This equation is 
    $$ y = i, x = 1, ix = y $$
    Imaginary pole resolve with Euler of equation. And, caltan idea enter
    with Imaginary rotation and rout of system.
    $$ i \sin 90^{\circ} = -1 $$
    $$ 1 \sin 90^{\circ} = i $$
    This number system express imaginary of reverse pole in circle function
    with dalanversian and alearletter of relate equation.
    
    These equation are concluded of being formula,
    $$ \pi^{e} \cong e^{\pi} $$
    This relation of neipia and pai number is mistery of Euler product 
    of integral manifold with anti-gravity and gravity equation stimulation.



    This equation mension to become with Heisenberg deprivate manifold.
    $$ [\Hat{\Hat{\varve}} |: \chi \to \Hat{\Hat{\alearletter}}]{}^{\ll
    {(p,q)}} $$
    And, This equation project with beta function of step equation,
    moreover, belong with element of gravity and antigravity 
    construct the double accelarate of Hat of quantum manifold.
    $$ = [\Hat{\Hat{\varve}} |: y \to \Hat{\Hat{\square}}]^{\ll \beta(
    \square, \alearletter)} $$
    D-group is Hashinate of madule entrance.
    $$ \varve |: \chi \to H\Psi^{D} \to \bigoplus {i\hbar} \otimes x_1
    \otimes x_2 \cdots $$

    Therefore, deltalic of Fastinate from Heisenberg of step function
    into deprivate manifold in Quantum operator, this equation fastcall from
    possibility and differential mantescue of equation.
  $$ \Phi(\delta(H\Psi))^{\nabla} = \scalebox{1}[2]{\variiint} \mathrm{cohom D\chi}(M) $$
  $$ = \lim_{n=0}^{\infty}({}_n P_{r} f(x)^{n} g(x)^{n-r})^{\nabla} $$

  After all, daybergence operator equal with step of deprivate function,
  and this function mechanize with dalancate of deprivate manifold.
  $$ \bigoplus{(i\hbar^{\nabla})}^{\oplus L} = {\Hat{\Hat{H\Psi}}^
  {i\hbar^{{i\hbar}^{\cdots}}}} $$
  $$ = H\Psi^{i\hbar^{'}} = i\hbar^{i\hbar^{'}} = e^{H\Psi\log H\Psi} $$

\end{document}


